\section{本章纲要：征服必须变成治理}

金朝的兴起不是辽亡、宋亡之间的一次快速换代。女真集团必须把东北联盟的军事能力，同燕云、华北的城市、户籍、粮道和旧政权官民接合起来；宋金联盟、北宋覆亡与绍兴和议则同时改变了金、南宋和仍在流动的居民的选择。问题不在于一方是否“取代”另一方，而在于征服所得的人口、土地和文书怎样被重新征收、保护或动员。

本章由此考察中都、河南、山东、东北和宋金边区的不同尺度。猛安谋克、州县、军户和市场并非整齐覆盖北方；迁徙、战事和地方协商使它们在各地有不同代价。官修史能建立战事与制度的脉络，碑刻、墓志、钱币、城址和地方材料则让接收与承受的局部过程可见，且提醒我们不要把行政类别当作固定族群或社会全貌。
